% 七帧教学图像组：显示顺序与解码顺序；宽约100mm。实际编码器可选择不同的参考结构。
\begin{tikzpicture}[x=1cm,y=1cm,font=\figurefont\small,
  frame/.style={draw=#1,fill=#1Fill,rounded corners=1.8pt,line width=0.8pt,
    minimum width=9mm,minimum height=8mm,inner sep=0.8mm},
  ref/.style={-{Stealth[length=1.6mm]},draw=FigureStroke,line width=0.7pt},
  bref/.style={ref,dashed,draw=FigureMuted},
  note/.style={text=FigureMuted,font=\figurefont\footnotesize}]
  \node[anchor=east] at (-0.6,2.2) {显示顺序};
  \foreach \i/\t/\c in {0/I/FigureBlue,1/B/FigureRed,2/B/FigureRed,3/P/FigureYellow,
                        4/B/FigureRed,5/B/FigureRed,6/P/FigureYellow}{
    \node[frame=\c] (d\i) at ({1.2*\i},2.2) {$\mathrm{\t}_{\i}$};
  }
  \draw[ref] (d0.north) to[bend left=35] node[above,note] {参考} (d3.north);
  \draw[ref] (d3.north) to[bend left=35] (d6.north);
  \draw[bref] (d0.south) to[bend right=40] (d1.south);
  \draw[bref] (d3.south) to[bend left=40] (d1.south);
  \draw[bref] (d3.south) to[bend right=40] (d4.south);
  \draw[bref] (d6.south) to[bend left=40] (d4.south);
  \node[anchor=east] at (-0.6,0) {解码顺序};
  \foreach \k/\i/\t/\c in {0/0/I/FigureBlue,1/3/P/FigureYellow,2/1/B/FigureRed,3/2/B/FigureRed,
                           4/6/P/FigureYellow,5/4/B/FigureRed,6/5/B/FigureRed}{
    \node[frame=\c] (e\k) at ({1.2*\k},0) {$\mathrm{\t}_{\i}$};
  }
  \foreach \k in {0,...,5}{
    \pgfmathtruncatemacro{\n}{\k+1}
    \draw[ref,draw=FigureMuted] (e\k.east) -- (e\n.west);
  }
\end{tikzpicture}
